\documentclass[11pt,letterpaper]{article}

% Required packages for professional RFP documents
\usepackage[margin=1in]{geometry}
\usepackage{graphicx}
\usepackage{hyperref}
\usepackage{fancyhdr}
\usepackage{booktabs}
\usepackage{enumitem}
\usepackage{xcolor}
\usepackage{titlesec}
\usepackage[english]{babel}
\usepackage{amsmath}
\usepackage{array}
\usepackage{longtable}
\usepackage{float}
\usepackage{pmboxdraw}
\usepackage{listings}
\usepackage{xcolor}
\usepackage[utf8]{inputenc}

% Page setup
\pagestyle{fancy}
\fancyhf{}
\fancyhead[L]{\textbf{RFP Response}}
\fancyhead[R]{\thepage}
\fancyfoot[C]{\textit{Confidential and Proprietary}}

% Hyperref setup
\hypersetup{
    colorlinks=true,
    linkcolor=blue,
    filecolor=magenta,
    urlcolor=cyan,
    pdftitle={RFP Response},
    pdfauthor={Your Company Name},
    pdfsubject={Request for Proposal Response},
    pdfkeywords={RFP, proposal, response}
}

% Custom colors
\definecolor{primaryblue}{RGB}{41, 128, 185}
\definecolor{darkgray}{RGB}{52, 73, 94}

% Listings configuration
\lstset{
    basicstyle=\ttfamily\footnotesize,
    commentstyle=\color{darkgray},
    keywordstyle=\color{primaryblue}\bfseries,
    breaklines=true,
    frame=single,
    numbers=left,
    numberstyle=\tiny\color{darkgray},
    showstringspaces=false,
    backgroundcolor=\color{gray!10},
    columns=flexible
}

% Section formatting
\titleformat{\section}
    {\normalfont\Large\bfseries\color{primaryblue}}
    {\thesection}{1em}{}

\titleformat{\subsection}
    {\normalfont\large\bfseries\color{darkgray}}
    {\thesubsection}{1em}{}

% Document information
\title{\textbf{\LARGE Daystar Roku App}\\
        \large [Project Name/RFP Number]}
\author{Alejandro Gutiérrez Franco\\
        Email: alejandrogfster@gmail.com\\
        Cel: +57 320 244 0076\\
        LinkedIn: \href{https://www.linkedin.com/in/fragualej}{linkedin.com/in/fragualej}}
\date{\today}

\begin{document}

\maketitle
\thispagestyle{empty}

\newpage
\tableofcontents
\newpage

% Development Setup & Foundation
\section{Development Setup \& Foundation}
\label{sec:development-setup}

\subsection{Version Control with Gitflow}
The project uses Git with Gitflow branching model for organized development:
\begin{itemize}
    \item \textbf{Main:} Production-ready code that's live or ready for release
    \item \textbf{Develop:} Integration branch containing latest completed features
    \item \textbf{Feature:} Individual feature development isolated from main codebase
    \item \textbf{Release:} Preparation branch for upcoming version deployment
    \item \textbf{Hotfix:} Emergency fixes for critical production issues
\end{itemize}

\subsection{Development Environment}
The project requires the following tools and technologies:
\begin{itemize}
    \item \textbf{BrightScript:} Core scripting language for Roku channel logic and functionality
    \item \textbf{SceneGraph:} UI framework providing XML-based component definitions with observer patterns
    \item \textbf{BrighterScript Compiler:} Check the entire project for syntax and program errors without needing to run on an actual Roku device
    \item \textbf{NodeJS:} Publish Roku projects to a Roku device and manage build automation tasks
\end{itemize}

\subsection{Folder Distribution}
\begin{center}
\begin{minipage}{0.6\textwidth}
\begin{verbatim}
├── src/
│   ├── manifest
│   ├── main.brs
│   ├── mainScene/
│   │   ├── mainScene.xml
│   │   └── mainScene.brs
│   ├── components/
│   │   ├── models/
│   │   ├── tasks/
│   │   ├── uiLogic/
│   │   ├── views/
│   │   └── widgets/
│   ├── images/
│   │   ├── backgrounds/
│   │   ├── icons/
│   │   └── splash/
│   └── source/
│       └── utils/
├── build/
│   └── compiled channel packages
└── pkg/
    └── distribution-ready files
\end{verbatim}
\end{minipage}
\end{center}

The \texttt{src/} directory contains all source code including the manifest file, main entry point, SceneGraph components, and assets. The \texttt{build/} directory stores compiled packages ready for testing, while \texttt{pkg/} holds final distribution files.

This structure follows Roku development best practices and supports the complete development-to-deployment workflow with modern tooling.

% Architecture & Design
\section{Architecture \& Design}
\label{sec:architecture-design}

\subsection{MVC Architecture for Roku Channels}
Roku channels follow a modified MVC pattern adapted for SceneGraph framework. The Model layer handles data management through repositories and API services. The View layer consists of SceneGraph XML components defining UI structure and layout. The Controller layer uses centralized BrightScript logic attached to MainScene for coordinating application flow and handling cross-component communication.

\subsection{Model Layer}
Located in \texttt{src/components/tasks/} for HTTP communication (httpTask) and \texttt{src/components/models/} for data structures. Models handle API communication, business logic, and data serialization. BaseModel classes provide common functionality for data validation and transformation. Repository classes manage data access and caching strategies for optimal performance on Roku devices.

\subsection{View Layer}
SceneGraph XML files in \texttt{src/components/views/} define the visual structure and layout. Each view component has corresponding BrightScript files for initialization and basic UI logic. Views communicate through observer patterns and field binding, maintaining separation from business logic.

\subsection{Controller Layer}
Located in \texttt{src/components/uiLogic/}, BrightScript scripts attached to MainScene manage application-wide coordination through observers. Controllers handle ItemFocused and ItemSelected events from UI components, manage ScreenStack for navigation flow between views, coordinate analytics tracking across the channel, manage TextToSpeech parsing for accessibility, and orchestrate the overall application state and transitions.

This architecture ensures maintainable code organization while leveraging Roku's SceneGraph strengths for TV interface development.


% UI Components Analysis
\section{UI Components Analysis}
\label{sec:ui-components}

\subsection{Base Components}
\label{sec:base-components}

Base components provide the foundational building blocks for the entire application. All other components extend or utilize these base components to ensure consistency and code reuse.

\begin{center}
\small
\begin{tabular}{@{}p{0.25\textwidth}p{0.70\textwidth}@{}}
\toprule
\textbf{Component} & \textbf{Description} \\
\midrule
BaseComponent & Core component with common props, state management, and lifecycle methods \\
BaseTile & Base tile component with image, title, and metadata layout \\
BaseView & Base view component with common layout, navigation, and error handling \\
BaseModel & Base data model providing common serialization, validation, and property mapping for API response objects \\
BaseTask & Base task component providing helpers, utils, and constants that extended tasks inherit \\
\bottomrule
\end{tabular}
\end{center}

\subsection{UI Components}
\begin{center}
\small
\begin{tabular}{@{}p{0.25\textwidth}p{0.70\textwidth}@{}}
\toprule
\textbf{Component} & \textbf{Description} \\
\midrule
AppHeader & Reusable header component with brand logo and optional search functionality \\
HeroBanner & Promotional section for both VOD and live content discovery \\
TileGrid & Container component for organizing tiles in grid layout \\
BaseTile & Base component with common tile functionality (image, title, metadata) \\
ChannelTile & Live channel display tile with favorite indicators and "LIVE" status \\
SearchResultTile & Search result display tile with standard media metadata \\
EpisodeTile & Episode display tile with duration and episode information \\
ContentDetails & Reusable metadata display component for shows, episodes, and live streams \\
PlaybackControls & Video controls including play buttons, resume, and play all options \\
SearchBox & Display area for search terms and query results \\
Keyboard & On-screen keyboard input system for TV/remote interfaces \\
LoadingIndicator & Loading spinner with text for async operations \\
\bottomrule
\end{tabular}
\end{center}

\subsection{Service Components}
Non-visual components that provide business logic, services, and utilities throughout the application:

\begin{center}
\small
\begin{tabular}{@{}p{0.25\textwidth}p{0.70\textwidth}@{}}
\toprule
\textbf{Component} & \textbf{Description} \\
\midrule
Repository & Data access layer coordinating HTTP operations and model transformations with caching strategies \\
TTSManager & Accessibility service providing audio narration for UI elements and content descriptions \\
Analytics & Tracking service for user interactions, screen views, and performance metrics \\
CacheManager & Data caching system for API responses and content optimization \\
AuthService & Authentication handling for user sessions and API authorization \\
ConfigManager & Application configuration management for settings and feature flags \\
Logger & Debugging and logging utility for development and production monitoring \\
DeepLinkHandler & URL routing service for external link navigation and content discovery \\
\bottomrule
\end{tabular}
\end{center}

\subsection{View Components}
The main UI Components table shows how each view utilizes the UI components:

\begin{center}
\small
\begin{longtable}{@{}p{0.25\textwidth}p{0.70\textwidth}@{}}
\toprule
\textbf{View} & \textbf{Components} \\
\midrule
HomeView & AppHeader, HeroBanner, TileGrid, ChannelTile \\
\midrule
LiveShowView & AppHeader, PlaybackControls, ContentDetails \\
\midrule
SearchView & SearchBox, Keyboard, TileGrid, SearchResultTile \\
\midrule
VodEpisodeView & AppHeader, PlaybackControls, ContentDetails \\
\midrule
VodShowView & AppHeader, ContentDetails, TileGrid, EpisodeTile \\
\midrule
LoadingView & AppHeader, LoadingIndicator \\
\midrule
ErrorView & Basic error display \\
\bottomrule
\caption{UI Components Breakdown}
\label{tab:ui-components}
\end{longtable}
\end{center}

\begin{center}
\small
\begin{tabular}{@{}lc@{}}
\toprule
\textbf{Component} & \textbf{Used in Views} \\
\midrule
AppHeader & 5 views \\
TileGrid & 3 views \\
BaseTile & 3 views \\
ChannelTile & 1 view \\
SearchResultTile & 1 view \\
EpisodeTile & 1 view \\
ContentDetails & 3 views \\
PlaybackControls & 2 views \\
LoadingIndicator & 1 view \\
HeroBanner & 1 view \\
SearchBox & 1 view \\
Keyboard & 1 view \\
\bottomrule
\end{tabular}
\end{center}

% Focus Handling
\section{Focus Handling}
\label{sec:focus-handling}

Focus management is a critical aspect of Roku application development, as users navigate through the interface using remote control directional keys. The Roku framework provides a robust focus handling system centered around the \texttt{onKeyEvent} function and native built-in navigation.

\subsection{The onKeyEvent Function}
\label{sec:onkeyevent}

The \texttt{onKeyEvent} function is the primary mechanism for handling remote control input in Roku applications. This function receives two parameters: \texttt{key} (string) containing the key name ("OK", "back", "up", "down", etc.) and \texttt{press} (boolean) indicating whether it's a press or release event. The function should return \texttt{true} to consume the event or \texttt{false} to allow it to bubble up to parent components. Always check \texttt{press = true} to handle only key press events and avoid duplicate actions.

\subsection{Native Component Navigation}
\label{sec:native-navigation}

Roku's native components like \texttt{RowList}, \texttt{MarkupList}, \texttt{MarkupGrid}, \texttt{TargetGroup}, \texttt{ButtonGroup}, and \texttt{LabelList} provide built-in navigation handling that automatically manages focus transitions without requiring custom \texttt{onKeyEvent} implementations. Simply assign focus to these components using \texttt{componentName.setFocus(true)} and they will handle all directional navigation, focus boundaries, visual indicators, and scroll behavior internally. This approach reduces code complexity while ensuring consistent navigation patterns optimized for TV interfaces.

% ScreenStack Management
\section{ScreenStack Management}
\label{sec:screenstack}

\subsection{ScreenStack Overview}
ScreenStack is custom navigation logic that needs to be developed to manage screen transitions and maintain navigation history. It handles the back button functionality and provides a consistent navigation experience across the channel.

\begin{lstlisting}[caption=ScreenStack Implementation Example]
' Push new screen to stack
screenData = { contentId: "12345", title: "Episode Details" }
screenStack.Push("EpisodeScreen", screenData)

' Handle back button (automatic)
' ScreenStack automatically pops current screen on back press

' Observe screen changes
screenStack.observeField("currentScreen", "onScreenChanged")

sub onScreenChanged()
  currentScreen = m.screenStack.currentScreen
  ' Handle screen transition logic
end sub
\end{lstlisting}

% API Services
\section{API Services}
\label{sec:api-services}

The API service layer provides structured HTTP communication and data management components for external service integration. The flow follows a circular pattern: View $\rightarrow$ HttpNode $\rightarrow$ Repository $\rightarrow$ HttpTask $\rightarrow$ HttpNode $\rightarrow$ View.

\begin{lstlisting}[caption=API Services Flow Implementation]
' View -> HttpNode (Request)
httpNode.request = { url: "/episodes", method: "GET" }

' HttpNode -> Repository
repository.callFunc("fetch", { httpNode: httpNode })

' Repository -> HttpTask
httpTask.requestData = {
  url: "https://api.example.com/episodes",
  httpNode: httpNode
}

' HttpTask processing (background)
while true
  msg = wait(0, m.port)
  msgType = type(msg)

  if msgType = "roSGNodeEvent"
    ' New request - escape chars and add headers
    request = msg.getData()
    request.url = escapeUrl(request.url)
    request.headers = addAuthHeaders(request.headers)
    processRequest(request)

  else if msgType = "roUrlEvent"
    ' Response received - parse JSON and create models
    response = ParseJson(msg.getString())
    model = CreateContentModel(response)

    ' HttpTask -> HttpNode (Response)
    targetHttpNode.response = model
  end if
end while

' HttpNode -> View (Response)
sub onHttpNodeResponse()
  m.homeView.content = m.httpNode.response.items
end sub
\end{lstlisting}

\subsection{HttpTask}
\label{sec:httptask}

Handles individual HTTP requests with URL, headers, and request body configuration. Manages request lifecycle and response parsing for REST API calls.

\subsection{HttpNode}
\label{sec:httpnode}

SceneGraph node wrapper for HTTP operations that bridges Roku's task-based architecture with HTTP functionality. Manages multiple HttpTasks and provides observable response handling.

\subsection{BaseModel}
\label{sec:basemodel}

Abstract data model providing common serialization, validation, and property mapping for API response objects. Standardizes data transformation between JSON and Roku objects.

\subsection{Repository}
\label{sec:repository}

Data access layer that coordinates HttpNode operations and BaseModel transformations. Implements caching, error handling, and provides clean API interface for business logic components.

% Certification Process
\section{Certification Process}
\label{sec:certification}

The Roku certification process ensures channels meet platform standards for functionality, performance, and user experience before publication to the Roku Channel Store.

\subsection{Certification Requirements}
TBD - Confirm criteria for Apps and Channels on Roku platform including:
\begin{itemize}
    \item Content guidelines and compliance standards
    \item Technical performance requirements and testing protocols
    \item User interface and accessibility compliance
    \item Security and privacy policy adherence
    \item Channel metadata and artwork specifications
    \item Deep linking and search integration requirements
\end{itemize}

\subsection{Submission Process}
The certification workflow involves channel package preparation, testing validation, and Roku review process for final approval and Channel Store publication.

% Appendices
\section{Appendices}
\label{sec:appendices}

Supporting documentation and additional materials.

\end{document}